\section{Grounded vs.\ Ungrounded Transformers}
\label{sec:grounded-ungrounded}

A grounded transformer and a standard language model may share the same
underlying architecture while differing substantially in how their
internal representations relate to externally defined concepts and
verification procedures.

Within the BP interpretation developed in this paper, both grounded and
ungrounded systems perform inference over implicit factor-graph-like
structures. The key distinction is whether the graph has a formally
declared interpretation.

\subsection*{The Distinction}

A \emph{grounded} transformer has:

\begin{itemize}
\item An explicit factor graph: tokens or latent states correspond to
declared propositions in a known knowledge base.

\item A finite concept space: the relevant propositions are bounded and
specified in advance.

\item A verifier: outputs can be checked against externally defined
correctness criteria.

\item Formal guarantees in restricted settings: on tree-structured
knowledge bases with constructive BP weights, exact posterior guarantees
follow from the companion theorem~\cite{coppola2026transformers}.
\end{itemize}

An \emph{ungrounded} transformer has:

\begin{itemize}
\item An implicit factor graph encoded in learned parameters rather than
explicitly declared structures.

\item No formally declared ontology tying internal states to a finite,
verified concept inventory.

\item No architecture-internal correctness guarantee analogous to the
grounded setting.

\item Reliability that is primarily empirical rather than formally
verified.
\end{itemize}

\subsection*{Hallucination in BP Terms}

Within this framework, hallucination can be interpreted as a mismatch
between model beliefs and externally verified posteriors.

For grounded systems, the comparison is well-defined:
\[
b(j) \neq P_{\mathrm{true}}(j),
\]
where $P_{\mathrm{true}}(j)$ is the correct posterior induced by the
knowledge base and observed evidence.

Ungrounded language models differ because they lack a formally declared
ontology and verifier. Their outputs may still correspond to meaningful
or true propositions in ordinary semantic practice, but the architecture
does not itself provide a mechanism guaranteeing correctness relative to
a fixed external knowledge base.

From this perspective, hallucination is not merely a numerical error in
the inference procedure. It is also a consequence of operating without a
fully specified and verifiable grounding structure.

\subsection*{Scaling and Grounding}

Scaling improves empirical performance. Larger models generally produce
more coherent outputs, achieve lower hallucination rates on benchmarks,
and exhibit richer internal representations.

However, scaling alone does not provide the same kind of formal
correctness guarantees available in grounded settings. Increasing model
capacity improves the quality of the learned implicit factor graph, but
it does not by itself create an explicit ontology, verifier, or formally
grounded concept space.

The distinction is therefore not between ``working'' and ``non-working''
systems. It is between empirical reliability and formally specified
correctness guarantees.

\subsection*{What Grounding Would Look Like}

A fully grounded transformer system would require several components:

\begin{enumerate}
\item A declared concept space: propositions or entities are explicitly
represented and bounded.

\item A declared relational structure: dependencies between concepts are
specified or verified.

\item A grounding procedure: natural language inputs are mapped into the
declared representational system.

\item A verification mechanism: outputs can be checked against the
underlying ontology or knowledge base.
\end{enumerate}

Retrieval-augmented generation (RAG) can be viewed as a partial form of
grounding because retrieved documents temporarily constrain the inference
space. More complete grounding would require persistent ontological and
verification structure rather than transient context retrieval alone.

\subsection*{The Practical Implication}

The practical distinction between grounded and ungrounded systems is the
kind of reliability guarantee they support.

For grounded systems in restricted settings, the constructive BP results
yield exact posterior guarantees. For general-purpose language models,
the strongest guarantees are typically empirical: benchmark performance,
human evaluation, calibration studies, or retrieval consistency.

The gap between these two forms of reliability --- formal verification
versus empirical robustness --- is the grounding gap.